% =========================================================================
\chapter{Likely questions and short answers}
\label{ch:questions}
% =========================================================================

\newcommand{\qa}[2]{\par\medskip\noindent\textbf{Q. #1}\par\noindent #2\par}

\qa{What is extreme value theory?}{The branch of statistics for the largest values. Its key result
(Pickands--Balkema--de Haan) says the excesses above a high threshold follow a generalised Pareto distribution
(GPD), whatever the original distribution.}

\qa{What does $\xi$ tell you?}{The shape of the tail. Negative (as here, about $-0.15$): the tail has an upper limit,
which is physical for temperature. Zero: exponential tail. Positive: heavy tail with no limit.}

\qa{Why not just oversample (SMOTE or copies)?}{Copies add no new information, and SMOTE only draws between existing
points, so neither reaches beyond the hottest day seen. The control shows random copies give 0.464 against 0.58--0.60
with EVT-guided generated days.}

\qa{How do you know the generated days are realistic?}{Every row is cleaned to obey physical rules and placed on a real
cell; the 8 gap columns are recomputed; and the test is done only on real days from 2016--2025, which the generators
never saw.}

\qa{Why the 99th percentile?}{It gives about 1\% extremes per cell and season, rare enough to be ``extreme'' but with
enough cases (29,604 in training) to learn from. It is computed separately per cell and season, so a hot day in Kinnaur
and in Una are judged on their own scales.}

\qa{Why 1981--2010 for thresholds?}{It is a standard 30-year WMO climate normal period, and it keeps the thresholds
fixed before the validation and test years.}

\qa{Why was the test set not split randomly?}{Days close in time are similar; a random split would leak information.
Validation and test are later, separate years.}

\qa{Is there leakage from tomorrow?}{No. All predictors use today and earlier; only the label looks at tomorrow.}

\qa{Why LightGBM?}{Best known model type for tables, fast on CPU, allows the physics (monotone) constraint. An MLP
scored the same on real data. One fixed classifier means differences come from the data.}

\qa{Why does CTGAN gain less from EVT?}{Without EVT it already produced many hot-tail days (plain score 0.554, the
highest), so there was less to fix.}

\qa{What is the 60\%?}{Tail positions are drawn half from the whole tail and half from its top fifth, so
$0.5\times0.2+0.5\times1=0.6$: 60\% of added days fall in the top 20\% of the tail.}

\qa{Why cap the severity?}{With $\xi<0$ the GPD has an end point $\sigma/|\xi|$; going past it is impossible. The cap
also stops anything above 1.5 times the hottest real excess.}

\qa{Why 529 cells?}{The ERA5-Land grid (0.1$^\circ$) over the box around Himachal Pradesh has $31\times37=1{,}147$
cells; 529 have their centre inside the state boundary.}

\qa{Why cap soil moisture at 0.8?}{Soil moisture is the fraction of the soil that is water and cannot exceed the
soil's saturation. In ERA5-Land's land model the largest saturation is 0.766 (organic soil), and our data never went
above 0.751. So 0.8 is a round limit just above anything real: it removes only impossible generated values and never
cuts a realistic one. It is our choice, not a constant; a per-cell limit would be stricter.}

\qa{Is the physics clean-up real physics?}{No, it is a consistency and range clean-up: rows are placed on a real
cell, values are cut to valid ranges, temperatures are ordered and the gap columns recomputed. It removes impossible
values but does not check a day against a weather model; humidity and the history columns are not recomputed.}

\qa{What would you do next?}{Longer lead times; let the GPD vary by cell or with elevation; check synthetic days with a
meteorologist; try the method in another region.}
